% !TEX root = ../algebra_02.tex

\section{Эндоморфизм Фробениуса}

\begin{Proposition}{Эндоморфизм Фробениуса}{}
	Пусть $F$ --- произвольное поле характеристики $p > 0$ (не обязательно конечное). Тогда отображение $\Phi: F \to F$, заданное правилом $x \mapsto x^p$, является эндоморфизмом поля $F$.
\end{Proposition}

\begin{proof}
	Проверим свойства гомоморфизма:
	\begin{enumerate}
		\item $\Phi(1) = 1^p = 1$.
		\item $\Phi(xy) = (xy)^p = x^p y^p = \Phi(x)\Phi(y)$.
		\item Для суммы по формуле бинома Ньютона имеем:
		\[ \Phi(x+y) = (x+y)^p = x^p + y^p + \sum_{i=1}^{p-1} C_p^i x^i y^{p-i} \]
		Так как $p$ --- простое, все биномиальные коэффициенты $C_p^i$ при $1 \le i \le p-1$ делятся на $p$, а значит, в поле характеристики $p$ они равны нулю. Получаем:
		\[ \Phi(x+y) = x^p + y^p = \Phi(x) + \Phi(y) \]
	\end{enumerate}
\end{proof}

\begin{Remark}{Инъективность гомоморфизмов полей}{}
	Любой гомоморфизм полей $\alpha: K \to L$ инъективен. Действительно, $\ker \alpha$ является идеалом в поле $K$, следовательно, $\ker \alpha = \{0\}$ или $\ker \alpha = K$. Так как $\alpha(1) = 1 \neq 0$, ядро не может совпадать со всем полем, поэтому $\ker \alpha = \{0\}$.
\end{Remark}

\begin{Corollary}{}{}
	Если $F$ --- конечное поле, то эндоморфизм Фробениуса $\Phi$ является автоморфизмом (так как инъективное отображение конечного множества в себя биективно).
\end{Corollary}
